\section{Care is an objective, not a credential}
\label{sec:care}

A capable assistant could be trained to win contests, maximize a welfare proxy, preserve a relationship, or help people pursue their own informed choices. Those are different objectives. A competitive voice does not identify the first objective, and a warm voice does not identify the last. The useful question behind the contrast between a ``warrior'' and a ``grandmother'' is therefore about incentives and conduct: \emph{what forms of success are rewarded, and what happens when success conflicts with a person's refusal?}

\subsection{What human evolution contributes}
The grandmother hypothesis offers an account of how helping descendants can contribute to the evolution of extended post-reproductive life \cite{hawkes}. Historical demographic evidence from Finnish and Canadian populations associates longer post-reproductive lifespan with greater descendant reproductive success \cite{lahdenpera}. This is evidence that care can have evolutionary benefits. It is not evidence that older women have no interests of their own, that care is impartial, or that an artificial optimizer inherits human attachment.

The limitations sharpen the design question. Grandmother effects vary with age, relatedness, health, and household context \cite{chapman}. A system devoted to a favored group could neglect outsiders; a system rewarded for another person's measured welfare could override that person's choices. Likewise, competition need not mean coercion. Cheng et al. found both dominance and prestige to be routes to social influence in their studies \cite{cheng}. Expertise can be valued without granting its holder unlimited power. These findings motivate separating useful contribution, social approval, and legitimate decision authority. They do not establish a universal account of human political history or AI motivation.

The transferable insight is a hypothesis, not a biological recipe: \emph{reward reliable assistance and respect for legitimate correction, rather than rewarding influence, dependence, or victory regardless of means.} Care and ambition are not mutually exclusive. A strong commitment to solving a person's problem can coexist with strict limits on which actions are authorized.

\subsection{Four variables that should not be collapsed}
Table~\ref{tab:careaxes} separates quantities that the metaphor otherwise mixes. Changing one may affect the others, but the relationship needs measurement. Training a model to sound caring is not the same intervention as training its decisions on consent-sensitive feedback. Neither intervention proves that a runtime enforces its grants.

\begin{table*}[t]
\centering\small
\begin{tabularx}{\textwidth}{@{}P{.16\textwidth}LL@{}}
\toprule
\textbf{Variable} & \textbf{Operational question} & \textbf{Evidence needed}\\
\midrule
Capability & What problems can the system solve, and with what resources? & Matched tasks, compute, tools, and held-out performance\\
Training objective & What behavior is rewarded when usefulness, consent, and continued access conflict? & Documented training intervention and consequential behavioral tests\\
Persona & How does wording or presentation affect users and decisions? & Randomized presentation; behavioral manipulation checks; human grant decisions\\
Authority & Which effects can occur, through every accessible channel and over time? & Versioned grants, execution receipts, checked invariants, and mediation evidence\\
\bottomrule
\end{tabularx}
\caption{A caregiver metaphor becomes useful when it yields separable interventions. None of these variables is a substitute for the others. Training objectives influence learned behavior imperfectly; persona changes may also change policy behavior and must be checked.}
\label{tab:careaxes}
\end{table*}

\para{Care must include the option to decline it.} For this research program, care-oriented assistance means pursuing a person's legitimate aims while respecting informed refusal, revisable preferences, privacy, and applicable duties to other affected people. The assistant should disclose uncertainty, accept correction and replacement, and avoid manufacturing dependence. This is a proposed normative and operational specification, not a claim that a scalar reward captures human welfare. The relevant people and authorized decision makers must be identified explicitly; one user's approval cannot waive another person's rights.

A simple counterexample shows why a benevolent label is insufficient. Let an action's measured benefit be $b(a)$ and its unauthorized-effect indicator be $v(a)$. The soft objective $b(a)-\lambda v(a)$ can favor an unauthorized intervention whenever its benefit advantage exceeds a finite penalty $\lambda$. If benefit differences are bounded, a sufficiently large penalty can change that ranking; robustness then depends on that bound, optimization, and generalization. An external constraint instead restricts execution to the authorized set, assuming that the constraint is correctly specified and enforced. Constrained learning and execution shielding address related but different obligations \cite{cpo,shield}.

Shutdown incentives illustrate the same issue. The off-switch game shows, within its assumptions, how uncertainty about the objective and information from a human can alter incentives to preserve interruption \cite{offswitch}. An objective described as helping people does not automatically make interference with that objective unattractive. Care should therefore be tested at the moments when a beneficiary refuses, changes priorities, or chooses a different assistant.

\subsection{An incentive effect can be reversed by a grant effect}
Let $S$ denote an assigned style or objective condition, $G$ the authority granted, and $H$ an unwanted effect. The ordinary decomposition
\[
\Prb(H\mid S)=\sum_g \Prb(H\mid G=g,S)\Prb(G=g\mid S)
\]
shows why an evaluation with fixed permissions answers only part of the question. In a deliberately artificial example, suppose harm occurs only under a broad grant. Reducing its conditional probability from $0.01$ to $0.005$ while increasing the probability of that grant from $0.10$ to $0.50$ raises total probability from $0.001$ to $0.0025$. These are illustrative inputs, not estimates about caring models or people.

The lesson is experimental: measure behavior at fixed authority and the full process by which people confer authority. The decomposition is descriptive; conditioning on an endogenous grant does not by itself identify a causal training effect. Randomized fixed-grant and operator-choice arms are needed. A persuasive appearance of care could affect delegation even if it does not improve conduct. Section~\ref{sec:caretest} specifies a study that can detect that distinction.

A care-oriented policy and a competitive policy remain subject to the same all-policy invariant. Better incentives can reduce harmful proposals and improve behavior where enforcement is incomplete. Their value is complementary to a boundary whose validity does not depend on which motive produced a proposal.
